\scanpage{6}
% 1–3
\uline{Conséquence}. Soit $k_0=A_0$, c'est un sous-corps de $k$. Soit $L(k_0)$ la catégorie des Alg. de Lie de dim finie sur $k_0$. Par Hausdorff on a un foncteur

% 4
\diagram{\begin{tikzcd}[ampersand replacement=\&,column sep=large]
L(k_0)\arrow[r,"\text{Hausdorff}"'] \& \mathrm{Grf}(k_0)\arrow[r,"\substack{\text{ext. de la}\\\text{base}}"'] \& \mathrm{Grf}(k)
\end{tikzcd}}

% 5–7
où ``Grf'' désigne la catégorie des groupes formels. Ceci dit, ce foncteur se factorise canoniquement par les $A$-groupes (formels) sur $k$.

% 8–10
De façon plus précise, soit $\overline W(\mathfrak g_0)$ le complété à l'origine \deleted{d'un}\inserted{de $\mathfrak g_0\otimes_{k_0}k$} pour \inserted{un} $k_0$-espace vectoriel \inserted{de rg fini} \deleted{$\unsure{k_0|\mathfrak g_0}$}. C'est un $A$-schéma formel.

% 11–14
Les formules de Hausdorff, si $\mathfrak g_0$ est une Alg de Lie, définissent une loi de groupe
\[
 \overline W(\mathfrak g_0)\times\overline W(\mathfrak g_0)\longrightarrow\overline W(\mathfrak g_0)
\]
\uline{Cette dernière est une $A$-loi} ; \deleted{Pour ceci, il suffit de vérifier \deleted{deux choses}\inserted{que le système de Pfaff}
1°) associé est donné par une $A$-loi \inserted{\unsure{un $A$-fibré}}} et si

% 15–20
$\mathfrak g_0\longrightarrow\mathfrak h_0$ est un hom d'algèbres de Lie sur $k_0$, l'hom de groupes formels $\overline W(\mathfrak g_0)\longrightarrow\overline W(\mathfrak h_0)$ est un $A$-hom. Utilisant le Th. d'Ado, on est ramené à prouver que si $\mathfrak g_0\subset\mathfrak{gl}(n,k_0)$, alors le Noyau de \inserted{ss/}gpe formel de $Gl(n,k_0)$ qu'il définit est un sous-$A$-schéma formel ....
% archive sheet number: 5
